\begin{pastexam}[label={px:ntu-113-dsa-10}]{臺大資工所 113 學年度 第 10 題}
（複選；每個選項單獨計分）

(5 points) Consider a sequence of numbers $(v_1,\ldots,v_n)$. Now we remove $k$ numbers from the
sequence so that we have $k+1$ non-empty segments of numbers. For example, consider $(2,3,8,1,4)$. If
we remove 8 then we have two segments $(2,3)$ and $(1,4)$.

Now we want to minimize the maximum sum of numbers of a segment. Let $m(1,n,k)$ be the answer, then
what is the correct recursion for $m$ when $1\le i<j\le n$ and $0<k\le(j-i)/2$ ?
\begin{enumerate}[label=\Alph*.]
\item $m(i,j,k)=\min_{x=i}^{j}\max(\sum_{y=x}^{j}v_y,m(i,x,k-1))$
\item $m(i,j,k)=\min_{x=i+1}^{j-1}\max(\sum_{y=x+1}^{j}v_y,m(i,x-1,k-1))$
\item $m(i,j,k)=\min_{x=i+1}^{j-1}\max(m(x+1,j,k-1),\sum_{y=i}^{x-1}v_y)$
\item $m(i,j,k)=\min_{v+w=k-1,v,w\ge0}\min_{x=i+1}^{j-1}(\max(m(i,x-1,v),m(x+1,j,w))$
\item $m(i,j,k)=\min_{v+w=k-1,v,w\ge0}\min_{x=i}^{j}(\max(m(i,x,v),m(x,j,w))$
\end{enumerate}
\end{pastexam}

\begin{pastsoln}
答案：B、C、D。

$m(i, j, k)$ 是「從 $v_i, \ldots, v_j$ 拿掉 $k$ 個數、留下 $k + 1$ 段非空的段，最大段和的最小值」。
拿掉的數不能在兩端，也不能相鄰（否則中間會有空段），所以 $j - i + 1 \ge 2k + 1$，就是題目的條件
$k \le (j - i)/2$。 不可能的子問題（例如 $k > (j - i)/2$）的值要當成 $+\infty$，取 $\min$ 時才不會被選到
（第 11 題 D 選項就在問這件事）；$k = 0$ 時 $m(i, j, 0) = \sum_{y=i}^{j} v_y$。 設計遞迴式就是問「最後
一步做了什麼選擇」（資結書第~18.1.5~節）。

D 的原式少一個右括號，照意思讀成 $\max$ 兩項取完再取 $\min$。
\begin{itemize}
\item A 錯。 把最後一段取成 $v_x..v_j$，前面是 $m(i, x, k-1)$，$v_x$ 同時算在兩邊，而且沒有任何數被拿
掉：$k - 1$ 次拿掉加上 $0$ 次，總共只拿掉 $k - 1$ 個。 例如 $(1, 1, 1)$、$k = 1$：正確答案是 $1$（拿掉中
間），A 算出 $\min_x \max(\sum_{y=x}^{3} v_y, m(1, x, 0))$，$x = 1, 2, 3$ 分別是 $\max(3, 1)$、$\max(2, 2)$、
$\max(1, 3)$，得到 $2$。
\item B 對。 看\textbf{最後一個}被拿掉的數 $v_x$。 它右邊 $v_{x+1}..v_j$ 是最後一段，要非空，$x \le j - 1$；
它左邊 $v_i..v_{x-1}$ 再拿掉 $k - 1$ 個，最好是 $m(i, x - 1, k - 1)$，也要非空，$x \ge i + 1$。 兩邊互不影響，
最大段和是兩者的較大值，再對 $x$ 取最小。 $x$ 太靠左、左邊放不下 $k - 1$ 個拿掉的數時，那一項是
$+\infty$，不影響結果。
\item C 對。 和 B 對稱：看\textbf{第一個}被拿掉的數 $v_x$，左邊 $v_i..v_{x-1}$ 是第一段，右邊
$v_{x+1}..v_j$ 再拿掉 $k - 1$ 個。
\item D 對。 隨便挑一個被拿掉的數 $v_x$，它把序列切成左右兩邊，左邊拿掉 $v$ 個、右邊拿掉 $w$ 個，
$v + w = k - 1$。 每個解都對應到這種切法（取它的任何一個被拿掉的數都行），兩邊各自取最好，所以式子
正確。 B 和 C 是 D 固定 $w = 0$ 或 $v = 0$ 的特例，所以 D 也包含了最佳解，只是多算了一些重複的情形。
\item E 錯。 $v_x$ 同時屬於兩邊，沒有被拿掉，兩邊合計只拿掉 $k - 1$ 個；而且 $v_x$ 在兩段裡各算一次。
用同一個例子 $(1, 1, 1)$、$k = 1$：$v = w = 0$，$x = 2$ 時是 $\max(m(1,2,0), m(2,3,0)) = 2$，$x = 1$ 時
$\max(1, 3) = 3$，最小是 $2$，不是 $1$。
\end{itemize}
\end{pastsoln}
